\documentclass[a4paper,10pt]{article}
\usepackage[utf8]{inputenc}
\usepackage[T1]{fontenc}
\usepackage[ngerman]{babel}
\usepackage{geometry}
% Adjusted margins for maximum density
\geometry{a4paper,left=1.2cm,right=1.2cm,top=1.2cm,bottom=1.2cm}
% Slightly expand the text height to force the page to fill
\addtolength{\textheight}{1.5cm}

\usepackage{lmodern}
\usepackage{xcolor}
\usepackage{hyperref}
\usepackage{titlesec}
\usepackage{enumitem}

\definecolor{titleblue}{RGB}{0,51,140}
\definecolor{TITLEBLUE}{RGB}{0,51,140}
\hypersetup{colorlinks=true, urlcolor=titleblue, linkcolor=titleblue, pdfborder={0 0 0}}

\pagestyle{empty}
\renewcommand{\baselinestretch}{0.92}  % Increased from 1.1 to 1.2

\titleformat{\section}
  {\color{titleblue}\Large\bfseries\uppercase}{}{0em}{}
  [\color{titleblue}\titlerule]
\titlespacing\section{0pt}{4pt}{2pt}  % Increased spacing

\setlist[itemize]{leftmargin=*,topsep=1pt,itemsep=1pt,parsep=0pt}  % Increased spacing
\setlist[itemize,2]{leftmargin=*,topsep=1pt,itemsep=1pt,parsep=0pt}

\begin{document}
\sloppy

\begin{center}
    {\Huge\bfseries\color{titleblue} AISSAOUI Ismail}\\[1pt]
    {\LARGE\bfseries Ingenieur für KI \& Data Science}\\\  % Increased spacing
    {    \href{mailto:i.aissaoui@esi-sba.dz}{i.aissaoui@esi-sba.dz} \quad|\quad
    +213 660 70 77 96 \quad|\quad Chlef, Algerien \\  % Increased spacing
    \href{https://linkedin.com/in/aissaoui-ismail-6a77a92a7}{LinkedIn} \quad|\quad
    \href{https://github.com/i-aissaoui}{GitHub} \quad|\quad
    \href{https://ismail-aissaoui.vercel.app}{Portfolio}
    }
\end{center}

\section{Berufliches Profil}
  % Make the text larger
Ingenieur für KI und Data Science, spezialisiert auf fortschrittliche maschinelle Lernarchitekturen, einschließlich Mamba-SSM, PINNs und GNNs. Expertise in der Entwicklung von extrem latenzarmen Edge-KI-Lösungen und skalierbaren Cloud-Modellen für Industrie und Medizin. Fundierte Kenntnisse im Hochleistungsrechnen (C++, Next.js) und moderner webbasierter 3D-Visualisierung.

\section{Technische Fähigkeiten}
\textbf{Sprachen:} Python, C++, Next.js, TypeScript, JavaScript, SQL \\
\textbf{KI/ML:} PyTorch, TensorFlow, Mamba, PINNs, GNNs, Transformers, DeepONet, KAN \\
\textbf{3D/Web:} React, Next.js, Three.js, React Three Fiber, FastAPI, TailwindCSS \\
\textbf{Data/MLOps:} MLflow, PyArrow, Docker, Linux, Git, PostgreSQL, Parquet \\
\textbf{Optimierung:} Post-Training Quantization (INT8), RL, Genetische Algorithmen

\section{Akademische Abschlussarbeiten}
\noindent\textbf{Ingenieurarbeit: Dual-Track Cloud-Edge Digitaler Zwilling für die vorausschauende Wartung von Gasturbinen} \hfill 2026 \\
\textit{Sonatrach (Abteilung für industrielle Überwachung)}
\begin{itemize}
    \item \textbf{Architektur:} Entwurf eines 3-stufigen Frameworks, das sicherheitskritische Edge-Modelle von Cloud-Analysen entkoppelt.
    \item \textbf{Optimierung:} Entwicklung von \textbf{TCNEdge} und Kompilierung von 1D-Dilated-Convolutions in INT8 ONNX C++, wodurch eine Latenz von 0,049 ms erreicht wurde.
    \item \textbf{Robustheit:} Entwicklung von \textbf{MambaTS} mit Missingness-Aware-Gating für eine ununterbrochene Inferenz bei Sensorausfällen.
    \item \textbf{Physikbasiert:} Implementierung von \textbf{KAN-PCNN} und \textbf{ST-KDFormer} zur Sicherstellung der thermodynamischen Konsistenz über PyTorch Autograd.
\end{itemize}

  % Increased spacing
\noindent\textbf{Masterarbeit: Hämodynamische digitale Zwillinge: Ein physikbasierter KI-Ansatz für die Prognostik in der Intensivmedizin} \hfill 2026 \\
\textit{Nationale Hochschule für Informatik (ESI-SBA)}
\begin{itemize}
    \item \textbf{Architektur:} Entwicklung von \textbf{Medical\_MambaTS} für die Inferenz am Krankenbett und des Dual-Branch \textbf{SST\_KODA\_MultiScale} für Cloud-Vorhersagen mit 100-Hz-Telemetrie.
    \item \textbf{Modellierung:} Implementierung eines 0D Windkessel \textbf{PI\_DeepONet} zur Durchsetzung kardiovaskulärer Fluiddynamik ohne patientenübergreifendes Neutraining.
    \item \textbf{Effizienz:} Entwicklung einer Zero-Copy-Pipeline mit PyArrow Parquet, um GPU-Engpässe zu beseitigen.
\end{itemize}

\section{Berufliche Projekte}
\noindent\textbf{Aura --- Enterprise KI-Handelsplattform} \hfill 2025 \\
\begin{itemize}
    \item \textbf{Vorhersagen:} Aufbau eines hybriden Systems mit LSTMs für hochpräzise Zeitreihenvorhersagen.
    \item \textbf{Strategie:} Implementierung von Deep Q-Learning (RL) zur Optimierung langfristiger Handelsstrategien.
    \item \textbf{Bereitstellung:} Entwicklung eines durchsatzstarken FastAPI-Backends mit WebSockets und einer Next.js-Benutzeroberfläche.
\end{itemize}

\vspace{6pt}
\noindent\textbf{Career Connect --- GNN-basiertes Rekrutierungssystem} \hfill 2024 -- 2025 \\
\begin{itemize}
    \item \textbf{Relationales Lernen:} Entwicklung einer skalierbaren Rekrutierungs-Engine unter Verwendung von GNNs zur Abbildung von Kandidaten-Job-Beziehungen.
    \item \textbf{Verarbeitung:} Nutzung von PyTorch Geometric und SBERT für tiefe semantische Embeddings.
    \item \textbf{Architektur:} Bereitstellung über FastAPI und PostgreSQL, wodurch die Übereinstimmungspräzision erheblich verbessert wurde.
\end{itemize}

\section{Bildungsweg}
\textbf{Nationale Hochschule für Informatik (ESI-SBA)} \hfill 2021 -- 2026 \\
Staatsingenieur und Masterabschluss in Künstlicher Intelligenz \& Data Science \\
\textit{Verteidigt am 28. Juni 2026}

\section{Zertifizierungen \& Sprachen}
\begin{itemize}[label=--]
    \item \textbf{DeepLearning.AI:} Machine Learning Specialization, Deep Learning Specialization, Natural Language Processing Specialization
    \item \textbf{Zero To Mastery (ZTM) Academy:} Complete PyTorch: Deep Learning and Artificial Intelligence, TensorFlow Developer Certificate 2024
    \item \textbf{Sprachen:} Arabisch (Muttersprache), Englisch (Verhandlungssicher), Französisch (Fließend)
\end{itemize}

\end{document}
